\section{Introduction}
\label{sec:introduction}

The usual comparison lower bound for sorting is a counting argument: a binary
decision tree distinguishing all $n!$ orders has a path of length at least
$\lceil\logtwo(n!)\rceil$. It does not immediately tell an adversary how to
find such a path efficiently. An adversary must answer each comparison as it
arrives, without knowing the sorter's future queries, and every answer must
remain consistent with distinct keys in a strict total order.

An unrestricted adversary can always choose the answer retaining more linear
extensions of the current partial order. The computational difficulty is that
exactly counting linear extensions is $\#\mathrm P$-complete
\citep{BrightwellWinkler1991}. This does not rule out an efficient optimal
adversary; it suggests looking for a different quantity to count.

Our adversary follows that route. It counts auxiliary integer assignments,
permits ties in those assignments, and recursively preserves the freedom to
resolve ties. Its organization extends the classical adversaries that move
items down an ordered tree. The main difference is that a move is charged for
the freedom it loses in its parent \emph{and credited} for the freedom it
creates in its child.

\begin{samepage}
\begin{theorem}[Main result]
\label{thm:main}
There is a single deterministic adversary which, initialized using only $n$,
answers comparisons consistently with a strict total order on $n$ distinct
keys, uses a fixed polynomial number of bit operations per answer, and forces
every correct comparison sorter to make
\[
                         (1-o(1))n\logtwo n
\]
comparisons.
\end{theorem}
\end{samepage}

Here ordinary sorting means comparison sorting of distinct, otherwise
unrestricted keys. The sorter can choose any pair adaptively. The adversary
does not receive a target accuracy or a description of the sorter. The
coefficient is optimal since merge sort uses at most $n\lceil\logtwo n\rceil$
comparisons. We do not claim an additive $O(n)$ gap from $\logtwo(n!)$, or
useful finite-size thresholds.

\subsection{The earlier adversaries}

\paragraph{Atallah and Kosaraju.}
Atallah and Kosaraju gave an explicit sorting adversary
\citep{AtallahKosaraju1981}. One useful way to view its effect is through
individual items: Richards and Vaidya describe how their technique forces,
for each $1\leq k<n$, $n-k$ items to participate in at least
$\logtwo(k+1)$ comparisons each \citep[Section~2]{RichardsVaidya1988}.
Each comparison has two participants. Taking, for example,
$k=\lfloor n/\logtwo n\rfloor$ gives the leading coefficient $1/2$ for
the total number of comparisons.

\paragraph{Borodin, Guibas, Lynch, and Yao.}
In their work on preprocessing for search, Borodin et al.\ gave a particularly
transparent bin-tree adversary \citep[Section~4]{BorodinEtAl1981}. Place all
items at the root of an infinite ordered binary tree. For a query on $x,y$:
\begin{enumerate}[label=(\roman*)]
\item If they occupy the same node, move $x$ to its left child and $y$ to
its right child; answer $x<y$.
\item If one's node properly contains the other in its subtree, move the
ancestor item to the child outside that subtree; answer according to the
resulting left-to-right order.
\item If their nodes already lie in separate subtrees, answer according to
their left-to-right order without moving either item.
\end{enumerate}
The spatial commitments preserve all earlier answers. Items at the same node
or at ancestor-related nodes remain publicly incomparable.

\paragraph{Brodal, Chaudhuri, and Radhakrishnan.}
Brodal et al.\ use this tree strategy in their work on maintaining rank
information and explicitly derive its sorting bound
\citep[Sections~4--4.1]{BrodalEtAl1996}. At termination the occupied nodes
are pairwise nonancestral and each contains one item. If their depths are
$d_1,\ldots,d_n$, the disjoint binary subtrees give
$\sum_i2^{-d_i}\leq1$. The arithmetic--geometric mean inequality yields
$\sum_i d_i\geq n\logtwo n$. A comparison increases the sum of depths by
at most two, giving $q\geq\frac12 n\logtwo n$.
This identifies the limitation of the depth accounting: a comparison can pay
for two downward steps.

Richards and Vaidya also give a binary-tree formulation and a weighted-poset
adversary \citep[Sections~3--4]{RichardsVaidya1988}. In the latter, item
weights are halved at comparisons while every chain has weight at most one.
The terminal inequality $\sum_x2^{-c(x)}\leq1$, with $c(x)$ the participation
count, gives the same coefficient. Kaplan, Zamir, and Zwick subsequently
improved the efficient deterministic guarantee to $37/64$
\citep{KaplanZamirZwick2019}.

\subsection{Our adversary in the moving-items framework}
\label{sec:framework}

We retain three features of the bin-tree view: items begin together, a move
commits an item to a more specific ordered region, and a comparison whose
endpoints still share a region is handled recursively there. We replace the
rule of separating the queried items immediately by retaining nearly half
of a counted family, then moving items when needed to simplify that family.

At a node, the regions are an ordered list of integer positions. A free item
can occupy an interval of these positions. Once assigned a fixed position,
it enters the recursive child associated with that position. Several items
can have the same parent position and remain unresolved in the child.
The integer positions are private auxiliary data, not key values.

This is an \emph{ordered multiway refinement} of the binary-tree framework.
A node's destinations can be drawn as leaves of a binary tree, but the
algorithm uses ordered lists. This expansion describes destinations; it does
not preserve the old ancestor-incomparability invariant. Free items can
already be publicly comparable, so we record their public order explicitly.

\begin{figure}[t]
\centering
\begin{tikzpicture}[>=Latex,font=\small,
  pool/.style={draw,rounded corners,align=center,minimum width=31mm,
               minimum height=10mm},
  slot/.style={draw,minimum width=8mm,minimum height=7mm}]
  \node[pool] (p) at (0,0) {free items\\$x,y,z,\ldots$};
  \node[slot] (a) at (-2.25,-1.35) {$1$};
  \node[slot] (b) at (-.75,-1.35) {$2$};
  \node[slot] (c) at (.75,-1.35) {$3$};
  \node[slot] (d) at (2.25,-1.35) {$4$};
  \draw[->] (p) -- (a);
  \draw[->] (p) -- (b);
  \draw[->] (p) -- (c);
  \draw[->] (p) -- (d);
  \node[pool] (child) at (-.75,-2.8) {child at position $2$\\$x,z,\ldots$};
  \draw[->] (b) -- (child) node[midway,right] {fix and admit};
  \node[align=left,text width=55mm] at (6,-1.4)
    {Equal auxiliary positions\\leave a recursive sorting problem.\\[4pt]
     An actual comparison is charged\\at only one node.};
\end{tikzpicture}
\caption{A move fixes an item's position in its parent and makes it free in
the corresponding child. The order between items sharing that position is
resolved recursively.}
\label{fig:recursion}
\end{figure}

The quantity replacing total depth is
\[
           \logtwo Z,\qquad Z=\prod_v Z_v,
\]
where $Z_v$ counts the node's admissible weakly order-preserving integer
assignments. Initially $Z=n^n$. At least one of the weak answers
$X_x\leq X_y$ and $X_y\leq X_x$ retains at least half the count:
the two events cover the family and overlap at ties. Small local
calculations approximate the needed probabilities.

To make those calculations reliable we keep each active graph sparse,
shallow, and controlled in its interval widths. Repairs can discard
assignments. But the $j$th item admitted to a child creates at least $j$ new
choices there. Repeated admissions therefore supply a factorial credit.
The proof shows that all net repair losses total $o(n\log n)$.
At completion the auxiliary count may still exceed one, but its logarithm
is only $O(n\log\log n)$. This suffices for the leading coefficient.

\paragraph{Organization and verification status.}
Section~\ref{sec:algorithm} gives the state, comparison procedure, and all
repair rules before their proofs. Sections~\ref{sec:legality}--\ref{sec:uniformity}
prove legality, local computation, and the lower bound. The complete local
approximation argument is in Appendix~\ref{app:local}.
The accompanying repository contains formal proofs of components of this
construction; it does not yet provide a complete Lean proof of
Theorem~\ref{thm:main}, including its controller and bit complexity.
The theorem here is supported by the mathematical argument below.

\paragraph{Notation.}
All logarithms have base two unless written $\ln$.
Put $T=n\logtwo n$ and interpret $0\logtwo0$ as zero.
The \emph{public order} $P$ is the transitive closure of actual comparison
answers; $x\leq_P y$ includes $x=y$. All algorithmic choices below use fixed
orders on labels, edges, paths, and blocks to break ties.
